
\documentclass[12pt]{article}

\usepackage[margin=1in]{geometry}
\usepackage{amsmath, amssymb}
\usepackage{changepage}
\usepackage{graphicx}
\usepackage{hyperref}

\setlength{\parindent}{0pt}
\setlength{\parskip}{10pt}

\title{Assessment -- Mathematical Induction}
\author{Serena Shao}
\date{November 9, 2026}

\begin{document}

\maketitle

\section*{\Huge\underline{Proof 1}}
\vspace{-0.5cm}
\section*{Question 8}
\vspace{-0.7cm}
\textbf{Problem--Induction Worksheet: Summations}

\vspace{0.3cm}

Use mathematical induction to prove that

\[
1\cdot 2^1+2\cdot 2^2+3\cdot 2^3
+\cdots+n\cdot 2^n
=(n-1)2^{n+1}+2
\]

for all positive integers $n$.

\textbf{Proof:}
We will proceed by mathematical induction on $n$.

\[
P(n):\quad
\sum_{i=1}^{n}i2^i=(n-1)2^{n+1}+2
\]

\subsection*{Base Case}

We first verify the statement for $n=0$.

Since the sum contains no terms when $n=0$,
the left-hand side is $0$.

\[
\begin{aligned}
\sum_{i=1}^{0}i2^i &= 0,\\
(0-1)2^{0+1}+2 &= -2+2=0.
\end{aligned}
\]

Thus, $P(0)$ is true, so the base case holds.

\subsection*{Inductive Step}

We need to show that $P(k)$ implies $P(k+1)$.

Assume $P(k)$ is true for some integer $k\geq 0$.

\[
P(k):\quad
\sum_{i=1}^{k}i2^i=(k-1)2^{k+1}+2
\]

We want to show that

\[
P(k+1):\quad
\sum_{i=1}^{k+1}i2^i=k2^{k+2}+2.
\]

Starting with the left-hand side of $P(k+1)$,

\[
\begin{aligned}
\sum_{i=1}^{k+1}i2^i
&=\sum_{i=1}^{k}i2^i+(k+1)2^{k+1}\\[6pt]
&=(k-1)2^{k+1}+2+(k+1)2^{k+1}\\[6pt]
&=(k-1+k+1)2^{k+1}+2\\[6pt]
&=2k\cdot 2^{k+1}+2\\[6pt]
&=k2^{k+2}+2.
\end{aligned}
\]

This is the statement $P(k+1)$.
Thus, $P(k)$ implies $P(k+1)$, so by induction,
$P(n)$ is true for all $n\in\mathbb{N}$.

\hfill $\square$

\vspace{0.5cm}

\newpage

\section*{\Huge\underline{Proof 2}}

\section*{Question 12}
\vspace{-0.7cm}
\textbf{Problems--Induction Worksheet: Miscellaneous}

\vspace{0.3cm}

Prove that any finite collection of squares of
various sizes can be cut into polygonal pieces
and rearranged to form a single square.

\textbf{Proof:}
We will proceed by mathematical induction on
the number of squares, $n$.

$P(n)$: Any collection of $n$ squares can be
cut into finitely many polygonal pieces and
rearranged to form a single square.

\subsection*{Base Case}

We begin with $n=1$, since the problem concerns
a nonempty collection of squares, so $n\geq 1$.

For $n=1$, the collection consists of exactly
one square, which already has the desired shape.
No cutting or rearrangement is necessary.

Thus, $P(1)$ is true.

\subsection*{Inductive Step}

We need to show that $P(k)$ implies $P(k+1)$.

Assume $P(k)$ is true for some integer $k\geq 1$.

Consider an arbitrary collection of $k+1$
squares with side lengths

\[
a_1,a_2,\ldots,a_k,a_{k+1}.
\]

By the inductive hypothesis, the first $k$
squares can be cut and rearranged into a
single square with side length $s$, where

\[
s^2=\sum_{i=1}^{k}a_i^2.
\]

We now need to combine this square with
the remaining square of side length $a_{k+1}$.

To justify this geometrically, consider two
arbitrary squares with side lengths $a$ and $b$.

Place four congruent right triangles, each
with legs $a$ and $b$, inside a larger square
of side length $a+b$, arranging them so that
their hypotenuses bound a central square.

\begin{center}
    \includegraphics[width=0.45\textwidth]{square_diagram.png}
\end{center}

By the Pythagorean theorem, the central
square has side length

\[
c=\sqrt{a^2+b^2}.
\]

The area of the larger square is

\[
(a+b)^2=a^2+2ab+b^2.
\]

The four triangles have total area

\[
4\left(\frac{ab}{2}\right)=2ab.
\]

Therefore, the area of the central square is

\[
(a+b)^2-2ab=a^2+b^2=c^2.
\]

The diagram shows that the square in the middle
has the same area as the two original squares
combined. However, having the same area does
not automatically mean we can rearrange the
squares into that shape. To do this, we can
cut up the two original squares into smaller
pieces and rearrange them so that they fit
together to form one square, without leaving
any gaps or overlapping.

This is important for our induction proof
because, after combining the first $k$ squares
into one, we can use the same cutting method
to combine that square with the remaining
square. Therefore, we can turn $k+1$ squares
into one square.

Applying this construction to the square
of side length $s$ and the square of side
length $a_{k+1}$, we obtain a single square
of side length $t$, where

\[
\begin{aligned}
t^2
&=s^2+a_{k+1}^2\\
&=\sum_{i=1}^{k}a_i^2+a_{k+1}^2\\
&=\sum_{i=1}^{k+1}a_i^2.
\end{aligned}
\]

This is the statement $P(k+1)$.
Thus, the original $k+1$ squares can be
cut into finitely many polygonal pieces
and rearranged into a single square.

This establishes $P(k+1)$.

Therefore, $P(k)$ implies $P(k+1)$.

\hfill $\square$

\vspace{0.5cm}

\textbf{Reviewed by: Pragya Khandelwal}


\vspace{0.5cm}

\textbf{Sources and AI Usage}

The diagram used in Proof 2 is not my original work.
I searched online for a diagram and found it here:
\url{https://www.geom.uiuc.edu/~demo5337/Group3/Choupei.html}

I used AI to help me learn how to use LaTeX,
understand its syntax and formatting, and navigate
Visual Studio Code (VS Code), including compiling
and troubleshooting my document.


\end{document}
